\documentclass[10pt,a4paper]{amsart}
\usepackage{amsmath,amssymb,amsthm,graphicx}
% For dealing with references we use the comment environment
\usepackage{verbatim}
\newenvironment{reference}{\comment}{\endcomment}
%\newenvironment{reference}{}{}
\newenvironment{slogan}{\comment}{\endcomment}
\newenvironment{history}{\comment}{\endcomment}
% For commutative diagrams we use Xy-pic
\usepackage[all]{xy}

% We use 2cell for 2-commutative diagrams.
\xyoption{2cell}
\UseAllTwocells
\theoremstyle{plain}
\newtheorem{theorem}[subsection]{Theorem}
\newtheorem{proposition}[subsection]{Proposition}
\newtheorem{lemma}[subsection]{Lemma}

\theoremstyle{definition}
\newtheorem{definition}[subsection]{Definition}
\newtheorem{example}[subsection]{Example}
\newtheorem{exercise}[subsection]{Exercise}
\newtheorem{situation}[subsection]{Situation}

\theoremstyle{remark}
\newtheorem{remark}[subsection]{Remark}
\newtheorem{remarks}[subsection]{Remarks}

\numberwithin{equation}{subsection}
% Macros
%
\def\lim{\mathop{\mathrm{lim}}\nolimits}
\def\colim{\mathop{\mathrm{colim}}\nolimits}
\def\Spec{\mathop{\mathrm{Spec}}}
\def\Hom{\mathop{\mathrm{Hom}}\nolimits}
\def\Ext{\mathop{\mathrm{Ext}}\nolimits}
\def\SheafHom{\mathop{\mathcal{H}\!\mathit{om}}\nolimits}
\def\SheafExt{\mathop{\mathcal{E}\!\mathit{xt}}\nolimits}
\def\Sch{\mathit{Sch}}
\def\Mor{\mathop{\mathrm{Mor}}\nolimits}
\def\Ob{\mathop{\mathrm{Ob}}\nolimits}
\def\Sh{\mathop{\mathit{Sh}}\nolimits}
\def\NL{\mathop{N\!L}\nolimits}
\def\CH{\mathop{\mathrm{CH}}\nolimits}
\def\proetale{{pro\text{-}\acute{e}tale}}
\def\etale{{\acute{e}tale}}
\def\QCoh{\mathit{QCoh}}
\def\Ker{\mathop{\mathrm{Ker}}}
\def\Im{\mathop{\mathrm{Im}}}
\def\Coker{\mathop{\mathrm{Coker}}}
\def\Coim{\mathop{\mathrm{Coim}}}

% Boxtimes
%
\DeclareMathSymbol{\boxtimes}{\mathbin}{AMSa}{"02}

%
% Macros for moduli stacks/spaces
%
\def\QCohstack{\mathcal{QC}\!\mathit{oh}}
\def\Cohstack{\mathcal{C}\!\mathit{oh}}
\def\Spacesstack{\mathcal{S}\!\mathit{paces}}
\def\Quotfunctor{\mathrm{Quot}}
\def\Hilbfunctor{\mathrm{Hilb}}
\def\Curvesstack{\mathcal{C}\!\mathit{urves}}
\def\Polarizedstack{\mathcal{P}\!\mathit{olarized}}
\def\Complexesstack{\mathcal{C}\!\mathit{omplexes}}
% \Pic is the operator that assigns to X its picard group, usage \Pic(X)
% \Picardstack_{X/B} denotes the Picard stack of X over B
% \Picardfunctor_{X/B} denotes the Picard functor of X over B
\def\Pic{\mathop{\mathrm{Pic}}\nolimits}
\def\Picardstack{\mathcal{P}\!\mathit{ic}}
\def\Picardfunctor{\mathrm{Pic}}
\def\Deformationcategory{\mathcal{D}\!\mathit{ef}}
\usepackage{cleveref}
\crefname{originallemma}{original lemma}{original lemmas}
\Crefname{originallemma}{Original lemma}{Original lemmas}
\makeatletter
\newlabel{lemma-bootstrap-diagonal}{{046K}{source}{}{sourcepointer-1}{}}
\newlabel{lemma-bootstrap-diagonal@cref}{{[originallemma][0][]046K}{[1][0][]source}{}{}{}}
\makeatother
\crefname{originallemma}{original lemma}{original lemmas}
\Crefname{originallemma}{Original lemma}{Original lemmas}
\makeatletter
\newlabel{lemma-surjective-flat-locally-finite-presentation}{{04S1}{source}{}{sourcepointer-2}{}}
\newlabel{lemma-surjective-flat-locally-finite-presentation@cref}{{[originallemma][0][]04S1}{[1][0][]source}{}{}{}}
\makeatother
\crefname{originalreference}{original reference}{original references}
\Crefname{originalreference}{Original reference}{Original references}
\makeatletter
\newlabel{spaces-groupoids-lemma-quotient-pre-equivalence-relation-restrict}{{046P}{source}{}{sourcepointer-3}{}}
\newlabel{spaces-groupoids-lemma-quotient-pre-equivalence-relation-restrict@cref}{{[originalreference][0][]046P}{[1][0][]source}{}{}{}}
\makeatother
\crefname{originalreference}{original reference}{original references}
\Crefname{originalreference}{Original reference}{Original references}
\makeatletter
\newlabel{spaces-more-groupoids-lemma-restrict-preserves-type}{{04RP}{source}{}{sourcepointer-4}{}}
\newlabel{spaces-more-groupoids-lemma-restrict-preserves-type@cref}{{[originalreference][0][]04RP}{[1][0][]source}{}{}{}}
\makeatother
\crefname{originalreference}{original reference}{original references}
\Crefname{originalreference}{Original reference}{Original references}
\makeatletter
\newlabel{spaces-morphisms-lemma-permanence-finite-type}{{0462}{source}{}{sourcepointer-5}{}}
\newlabel{spaces-morphisms-lemma-permanence-finite-type@cref}{{[originalreference][0][]0462}{[1][0][]source}{}{}{}}
\makeatother
\crefname{originalreference}{original reference}{original references}
\Crefname{originalreference}{Original reference}{Original references}
\makeatletter
\newlabel{spaces-morphisms-lemma-monomorphism-loc-finite-type-loc-quasi-finite}{{0463}{source}{}{sourcepointer-6}{}}
\newlabel{spaces-morphisms-lemma-monomorphism-loc-finite-type-loc-quasi-finite@cref}{{[originalreference][0][]0463}{[1][0][]source}{}{}{}}
\makeatother
\crefname{originalreference}{original reference}{original references}
\Crefname{originalreference}{Original reference}{Original references}
\makeatletter
\newlabel{spaces-morphisms-proposition-locally-quasi-finite-separated-over-scheme}{{03XX}{source}{}{sourcepointer-7}{}}
\newlabel{spaces-morphisms-proposition-locally-quasi-finite-separated-over-scheme@cref}{{[originalreference][0][]03XX}{[1][0][]source}{}{}{}}
\makeatother
\crefname{originallemma}{original lemma}{original lemmas}
\Crefname{originallemma}{Original lemma}{Original lemmas}
\makeatletter
\newlabel{lemma-slice-equivalence-relation}{{0489}{source}{}{sourcepointer-8}{}}
\newlabel{lemma-slice-equivalence-relation@cref}{{[originallemma][0][]0489}{[1][0][]source}{}{}{}}
\makeatother
\crefname{originallemma}{original lemma}{original lemmas}
\Crefname{originallemma}{Original lemma}{Original lemmas}
\makeatletter
\newlabel{lemma-better-finding-opens}{{046M}{source}{}{sourcepointer-9}{}}
\newlabel{lemma-better-finding-opens@cref}{{[originallemma][0][]046M}{[1][0][]source}{}{}{}}
\makeatother
\crefname{originallemma}{original lemma}{original lemmas}
\Crefname{originallemma}{Original lemma}{Original lemmas}
\makeatletter
\newlabel{spaces-lemma-finding-opens}{{02WU}{source}{}{sourcepointer-10}{}}
\newlabel{spaces-lemma-finding-opens@cref}{{[originallemma][0][]02WU}{[1][0][]source}{}{}{}}
\makeatother
\crefname{originallemma}{original lemma}{original lemmas}
\Crefname{originallemma}{Original lemma}{Original lemmas}
\makeatletter
\newlabel{spaces-lemma-coproduct-algebraic-spaces}{{02WQ}{source}{}{sourcepointer-11}{}}
\newlabel{spaces-lemma-coproduct-algebraic-spaces@cref}{{[originallemma][0][]02WQ}{[1][0][]source}{}{}{}}
\makeatother
\crefname{originallemma}{original lemma}{original lemmas}
\Crefname{originallemma}{Original lemma}{Original lemmas}
\makeatletter
\newlabel{spaces-lemma-glueing-algebraic-spaces}{{02WR}{source}{}{sourcepointer-12}{}}
\newlabel{spaces-lemma-glueing-algebraic-spaces@cref}{{[originallemma][0][]02WR}{[1][0][]source}{}{}{}}
\makeatother
\crefname{originalreference}{original reference}{original references}
\Crefname{originalreference}{Original reference}{Original references}
\makeatletter
\newlabel{spaces-more-groupoids-lemma-quasi-splitting-affine-scheme}{{04S0}{source}{}{sourcepointer-13}{}}
\newlabel{spaces-more-groupoids-lemma-quasi-splitting-affine-scheme@cref}{{[originalreference][0][]04S0}{[1][0][]source}{}{}{}}
\makeatother
\crefname{originallemma}{original lemma}{original lemmas}
\Crefname{originallemma}{Original lemma}{Original lemmas}
\makeatletter
\newlabel{lemma-divide-subgroupoid}{{04S4}{source}{}{sourcepointer-14}{}}
\newlabel{lemma-divide-subgroupoid@cref}{{[originallemma][0][]04S4}{[1][0][]source}{}{}{}}
\makeatother
\crefname{originallemma}{original lemma}{original lemmas}
\Crefname{originallemma}{Original lemma}{Original lemmas}
\makeatletter
\newlabel{descent-lemma-flat-fpqc-local-source}{{036K}{source}{}{sourcepointer-15}{}}
\newlabel{descent-lemma-flat-fpqc-local-source@cref}{{[originallemma][0][]036K}{[1][0][]source}{}{}{}}
\makeatother
\crefname{originallemma}{original lemma}{original lemmas}
\Crefname{originallemma}{Original lemma}{Original lemmas}
\makeatletter
\newlabel{descent-lemma-locally-finite-presentation-fppf-local-source}{{036N}{source}{}{sourcepointer-16}{}}
\newlabel{descent-lemma-locally-finite-presentation-fppf-local-source@cref}{{[originallemma][0][]036N}{[1][0][]source}{}{}{}}
\makeatother
\crefname{originallemma}{original lemma}{original lemmas}
\Crefname{originallemma}{Original lemma}{Original lemmas}
\makeatletter
\newlabel{morphisms-lemma-quasi-finite-local-source}{{0GWS}{source}{}{sourcepointer-17}{}}
\newlabel{morphisms-lemma-quasi-finite-local-source@cref}{{[originallemma][0][]0GWS}{[1][0][]source}{}{}{}}
\makeatother
\crefname{originallemma}{original lemma}{original lemmas}
\Crefname{originallemma}{Original lemma}{Original lemmas}
\makeatletter
\newlabel{descent-lemma-descending-property-flat}{{02L2}{source}{}{sourcepointer-18}{}}
\newlabel{descent-lemma-descending-property-flat@cref}{{[originallemma][0][]02L2}{[1][0][]source}{}{}{}}
\makeatother
\crefname{originallemma}{original lemma}{original lemmas}
\Crefname{originallemma}{Original lemma}{Original lemmas}
\makeatletter
\newlabel{descent-lemma-descending-property-locally-finite-presentation}{{02KY}{source}{}{sourcepointer-19}{}}
\newlabel{descent-lemma-descending-property-locally-finite-presentation@cref}{{[originallemma][0][]02KY}{[1][0][]source}{}{}{}}
\makeatother
\crefname{originallemma}{original lemma}{original lemmas}
\Crefname{originallemma}{Original lemma}{Original lemmas}
\makeatletter
\newlabel{descent-lemma-descending-property-quasi-finite}{{02VI}{source}{}{sourcepointer-20}{}}
\newlabel{descent-lemma-descending-property-quasi-finite@cref}{{[originallemma][0][]02VI}{[1][0][]source}{}{}{}}
\makeatother
\crefname{originallemma}{original lemma}{original lemmas}
\Crefname{originallemma}{Original lemma}{Original lemmas}
\makeatletter
\newlabel{descent-lemma-descending-property-open-immersion}{{02L3}{source}{}{sourcepointer-21}{}}
\newlabel{descent-lemma-descending-property-open-immersion@cref}{{[originallemma][0][]02L3}{[1][0][]source}{}{}{}}
\makeatother
\crefname{originallemma}{original lemma}{original lemmas}
\Crefname{originallemma}{Original lemma}{Original lemmas}
\makeatletter
\newlabel{groupoids-lemma-quotient-pre-equivalence-relation-restrict}{{045Z}{source}{}{sourcepointer-22}{}}
\newlabel{groupoids-lemma-quotient-pre-equivalence-relation-restrict@cref}{{[originallemma][0][]045Z}{[1][0][]source}{}{}{}}
\makeatother
\crefname{originalreference}{original reference}{original references}
\Crefname{originalreference}{Original reference}{Original references}
\makeatletter
\newlabel{more-groupoids-lemma-sheaf-differentials}{{04R9}{source}{}{sourcepointer-23}{}}
\newlabel{more-groupoids-lemma-sheaf-differentials@cref}{{[originalreference][0][]04R9}{[1][0][]source}{}{}{}}
\makeatother
\crefname{originallemma}{original lemma}{original lemmas}
\Crefname{originallemma}{Original lemma}{Original lemmas}
\makeatletter
\newlabel{morphisms-lemma-unramified-omega-zero}{{02G5}{source}{}{sourcepointer-24}{}}
\newlabel{morphisms-lemma-unramified-omega-zero@cref}{{[originallemma][0][]02G5}{[1][0][]source}{}{}{}}
\makeatother
\crefname{originallemma}{original lemma}{original lemmas}
\Crefname{originallemma}{Original lemma}{Original lemmas}
\makeatletter
\newlabel{morphisms-lemma-flat-unramified-etale}{{02GV}{source}{}{sourcepointer-25}{}}
\newlabel{morphisms-lemma-flat-unramified-etale@cref}{{[originallemma][0][]02GV}{[1][0][]source}{}{}{}}
\makeatother
\crefname{originaltheorem}{original theorem}{original theorems}
\Crefname{originaltheorem}{Original theorem}{Original theorems}
\makeatletter
\newlabel{spaces-theorem-presentation}{{02WW}{source}{}{sourcepointer-26}{}}
\newlabel{spaces-theorem-presentation@cref}{{[originaltheorem][0][]02WW}{[1][0][]source}{}{}{}}
\makeatother
\title{S001-theorem-final-bootstrap}
\author{The Stacks Project Authors}
\date{stacks-project:a04446e57ec1fbc252a871afcec7752fb2807b14; GNU Free Documentation License 1.2 (as published in the Stacks Project COPYING)}
\begin{document}
\maketitle
\noindent\textit{Editorial note.} All mathematical statement and proof text is reproduced verbatim from the author's chapter bootstrap.tex of the Stacks Project, fixed at commit a04446e57ec1fbc252a871afcec7752fb2807b14. The pinned preamble is supplied as its own file; the document wrapper class, the xr-hyper wiring and the \textbackslash{}externaldocument list are omitted because the assembler supplies the class.\par\medskip
\noindent\textit{Editorial note.} Numbering inside this standalone excerpt is editorial, not the full Stacks book print numbering.\par\medskip
\noindent\textit{Editorial note.} The candidate was previously judged H2 by the initial-corpus stage; that judgement and its reason are reused verbatim and were not re-derived here. Original candidate key: S001-theorem-final-bootstrap.\par\medskip
\noindent\textit{Editorial note.} Reference labels in the proof resolve to navigation pointers to the original author source; these pointers do not reproduce proofs of the prerequisites.\par\medskip
\noindent\textit{Allowed standard prerequisites (editorial).}\begin{itemize}
\item Prerequisites named by the author inside the counted span. The counted span is reproduced verbatim, including the author's own citations. Its prerequisites are the earlier results of the same fixed Stacks Project version that the author cites there and in the chapter's lead-in.
\end{itemize}\medskip
\ifdefined\hypertarget\hypertarget{sourcepointer-1}{}\fi
\noindent\textit{Source reference.} Original lemma 046K: sources/bootstrap.tex, source lines 747--747. Navigation pointer used in the counted span; the span contains the pointer only and the cited result is not reproduced inside it. Registered as a non-load-bearing navigation stub in the established form for Stacks material. No bibliography citation is involved.\par\medskip
\ifdefined\hypertarget\hypertarget{sourcepointer-2}{}\fi
\noindent\textit{Source reference.} Original lemma 04S1: sources/bootstrap.tex, source lines 591--591. Navigation pointer used in the counted span; the span contains the pointer only and the cited result is not reproduced inside it. Registered as a non-load-bearing navigation stub in the established form for Stacks material. No bibliography citation is involved.\par\medskip
\ifdefined\hypertarget\hypertarget{sourcepointer-3}{}\fi
\noindent\textit{Source reference.} Original reference 046P: sources/spaces-groupoids.tex, source lines 2148--2148. Navigation pointer used in the counted span; the span contains the pointer only and the cited result is not reproduced inside it. Registered as a non-load-bearing navigation stub in the established form for Stacks material. No bibliography citation is involved.\par\medskip
\ifdefined\hypertarget\hypertarget{sourcepointer-4}{}\fi
\noindent\textit{Source reference.} Original reference 04RP: sources/spaces-more-groupoids.tex, source lines 585--585. Navigation pointer used in the counted span; the span contains the pointer only and the cited result is not reproduced inside it. Registered as a non-load-bearing navigation stub in the established form for Stacks material. No bibliography citation is involved.\par\medskip
\ifdefined\hypertarget\hypertarget{sourcepointer-5}{}\fi
\noindent\textit{Source reference.} Original reference 0462: sources/spaces-morphisms.tex, source lines 4489--4489. Navigation pointer used in the counted span; the span contains the pointer only and the cited result is not reproduced inside it. Registered as a non-load-bearing navigation stub in the established form for Stacks material. No bibliography citation is involved.\par\medskip
\ifdefined\hypertarget\hypertarget{sourcepointer-6}{}\fi
\noindent\textit{Source reference.} Original reference 0463: sources/spaces-morphisms.tex, source lines 5361--5361. Navigation pointer used in the counted span; the span contains the pointer only and the cited result is not reproduced inside it. Registered as a non-load-bearing navigation stub in the established form for Stacks material. No bibliography citation is involved.\par\medskip
\ifdefined\hypertarget\hypertarget{sourcepointer-7}{}\fi
\noindent\textit{Source reference.} Original reference 03XX: sources/spaces-morphisms.tex, source lines 10590--10590. Navigation pointer used in the counted span; the span contains the pointer only and the cited result is not reproduced inside it. Registered as a non-load-bearing navigation stub in the established form for Stacks material. No bibliography citation is involved.\par\medskip
\ifdefined\hypertarget\hypertarget{sourcepointer-8}{}\fi
\noindent\textit{Source reference.} Original lemma 0489: sources/bootstrap.tex, source lines 1121--1121. Navigation pointer used in the counted span; the span contains the pointer only and the cited result is not reproduced inside it. Registered as a non-load-bearing navigation stub in the established form for Stacks material. No bibliography citation is involved.\par\medskip
\ifdefined\hypertarget\hypertarget{sourcepointer-9}{}\fi
\noindent\textit{Source reference.} Original lemma 046M: sources/bootstrap.tex, source lines 990--990. Navigation pointer used in the counted span; the span contains the pointer only and the cited result is not reproduced inside it. Registered as a non-load-bearing navigation stub in the established form for Stacks material. No bibliography citation is involved.\par\medskip
\ifdefined\hypertarget\hypertarget{sourcepointer-10}{}\fi
\noindent\textit{Source reference.} Original lemma 02WU: sources/spaces.tex, source lines 1310--1310. Navigation pointer used in the counted span; the span contains the pointer only and the cited result is not reproduced inside it. Registered as a non-load-bearing navigation stub in the established form for Stacks material. No bibliography citation is involved.\par\medskip
\ifdefined\hypertarget\hypertarget{sourcepointer-11}{}\fi
\noindent\textit{Source reference.} Original lemma 02WQ: sources/spaces.tex, source lines 1060--1060. Navigation pointer used in the counted span; the span contains the pointer only and the cited result is not reproduced inside it. Registered as a non-load-bearing navigation stub in the established form for Stacks material. No bibliography citation is involved.\par\medskip
\ifdefined\hypertarget\hypertarget{sourcepointer-12}{}\fi
\noindent\textit{Source reference.} Original lemma 02WR: sources/spaces.tex, source lines 1082--1082. Navigation pointer used in the counted span; the span contains the pointer only and the cited result is not reproduced inside it. Registered as a non-load-bearing navigation stub in the established form for Stacks material. No bibliography citation is involved.\par\medskip
\ifdefined\hypertarget\hypertarget{sourcepointer-13}{}\fi
\noindent\textit{Source reference.} Original reference 04S0: sources/spaces-more-groupoids.tex, source lines 2610--2610. Navigation pointer used in the counted span; the span contains the pointer only and the cited result is not reproduced inside it. Registered as a non-load-bearing navigation stub in the established form for Stacks material. No bibliography citation is involved.\par\medskip
\ifdefined\hypertarget\hypertarget{sourcepointer-14}{}\fi
\noindent\textit{Source reference.} Original lemma 04S4: sources/bootstrap.tex, source lines 1301--1301. Navigation pointer used in the counted span; the span contains the pointer only and the cited result is not reproduced inside it. Registered as a non-load-bearing navigation stub in the established form for Stacks material. No bibliography citation is involved.\par\medskip
\ifdefined\hypertarget\hypertarget{sourcepointer-15}{}\fi
\noindent\textit{Source reference.} Original lemma 036K: sources/descent.tex, source lines 6765--6765. Navigation pointer used in the counted span; the span contains the pointer only and the cited result is not reproduced inside it. Registered as a non-load-bearing navigation stub in the established form for Stacks material. No bibliography citation is involved.\par\medskip
\ifdefined\hypertarget\hypertarget{sourcepointer-16}{}\fi
\noindent\textit{Source reference.} Original lemma 036N: sources/descent.tex, source lines 6811--6811. Navigation pointer used in the counted span; the span contains the pointer only and the cited result is not reproduced inside it. Registered as a non-load-bearing navigation stub in the established form for Stacks material. No bibliography citation is involved.\par\medskip
\ifdefined\hypertarget\hypertarget{sourcepointer-17}{}\fi
\noindent\textit{Source reference.} Original lemma 0GWS: sources/morphisms.tex, source lines 4021--4021. Navigation pointer used in the counted span; the span contains the pointer only and the cited result is not reproduced inside it. Registered as a non-load-bearing navigation stub in the established form for Stacks material. No bibliography citation is involved.\par\medskip
\ifdefined\hypertarget\hypertarget{sourcepointer-18}{}\fi
\noindent\textit{Source reference.} Original lemma 02L2: sources/descent.tex, source lines 5833--5833. Navigation pointer used in the counted span; the span contains the pointer only and the cited result is not reproduced inside it. Registered as a non-load-bearing navigation stub in the established form for Stacks material. No bibliography citation is involved.\par\medskip
\ifdefined\hypertarget\hypertarget{sourcepointer-19}{}\fi
\noindent\textit{Source reference.} Original lemma 02KY: sources/descent.tex, source lines 5767--5767. Navigation pointer used in the counted span; the span contains the pointer only and the cited result is not reproduced inside it. Registered as a non-load-bearing navigation stub in the established form for Stacks material. No bibliography citation is involved.\par\medskip
\ifdefined\hypertarget\hypertarget{sourcepointer-20}{}\fi
\noindent\textit{Source reference.} Original lemma 02VI: sources/descent.tex, source lines 6117--6117. Navigation pointer used in the counted span; the span contains the pointer only and the cited result is not reproduced inside it. Registered as a non-load-bearing navigation stub in the established form for Stacks material. No bibliography citation is involved.\par\medskip
\ifdefined\hypertarget\hypertarget{sourcepointer-21}{}\fi
\noindent\textit{Source reference.} Original lemma 02L3: sources/descent.tex, source lines 5862--5862. Navigation pointer used in the counted span; the span contains the pointer only and the cited result is not reproduced inside it. Registered as a non-load-bearing navigation stub in the established form for Stacks material. No bibliography citation is involved.\par\medskip
\ifdefined\hypertarget\hypertarget{sourcepointer-22}{}\fi
\noindent\textit{Source reference.} Original lemma 045Z: sources/groupoids.tex, source lines 3819--3819. Navigation pointer used in the counted span; the span contains the pointer only and the cited result is not reproduced inside it. Registered as a non-load-bearing navigation stub in the established form for Stacks material. No bibliography citation is involved.\par\medskip
\ifdefined\hypertarget\hypertarget{sourcepointer-23}{}\fi
\noindent\textit{Source reference.} Original reference 04R9: sources/more-groupoids.tex, source lines 117--117. Navigation pointer used in the counted span; the span contains the pointer only and the cited result is not reproduced inside it. Registered as a non-load-bearing navigation stub in the established form for Stacks material. No bibliography citation is involved.\par\medskip
\ifdefined\hypertarget\hypertarget{sourcepointer-24}{}\fi
\noindent\textit{Source reference.} Original lemma 02G5: sources/morphisms.tex, source lines 7646--7646. Navigation pointer used in the counted span; the span contains the pointer only and the cited result is not reproduced inside it. Registered as a non-load-bearing navigation stub in the established form for Stacks material. No bibliography citation is involved.\par\medskip
\ifdefined\hypertarget\hypertarget{sourcepointer-25}{}\fi
\noindent\textit{Source reference.} Original lemma 02GV: sources/morphisms.tex, source lines 8405--8405. Navigation pointer used in the counted span; the span contains the pointer only and the cited result is not reproduced inside it. Registered as a non-load-bearing navigation stub in the established form for Stacks material. No bibliography citation is involved.\par\medskip
\ifdefined\hypertarget\hypertarget{sourcepointer-26}{}\fi
\noindent\textit{Source reference.} Original theorem 02WW: sources/spaces.tex, source lines 1603--1603. Navigation pointer used in the counted span; the span contains the pointer only and the cited result is not reproduced inside it. Registered as a non-load-bearing navigation stub in the established form for Stacks material. No bibliography citation is involved.\par\medskip
\begin{theorem}
\label{theorem-final-bootstrap}
Let $S$ be a scheme.
Let $F : (\Sch/S)_{fppf}^{opp} \to \textit{Sets}$ be a functor.
Any one of the following conditions implies that $F$ is an algebraic space:
\begin{enumerate}
\item $F = U/R$ where $(U, R, s, t, c)$ is a groupoid in algebraic spaces
over $S$ such that $s, t$ are flat and locally of finite presentation, and
$j = (t, s) : R \to U \times_S U$ is an equivalence relation,
\item $F = U/R$ where $(U, R, s, t, c)$ is a groupoid scheme
over $S$ such that $s, t$ are flat and locally of finite presentation, and
$j = (t, s) : R \to U \times_S U$ is an equivalence relation,
\item $F$ is a sheaf and there exists an algebraic space $U$ and a morphism
$U \to F$ which is representable by algebraic spaces,
surjective, flat and locally of finite presentation,
\item $F$ is a sheaf and there exists a scheme $U$ and a morphism
$U \to F$ which is representable by algebraic spaces or schemes,
surjective, flat and locally of finite presentation,
\item $F$ is a sheaf, $\Delta_F$ is representable by algebraic spaces,
and there exists an algebraic space $U$ and a morphism $U \to F$ which is
surjective, flat, and locally of finite presentation, or
\item $F$ is a sheaf, $\Delta_F$ is representable,
and there exists a scheme $U$ and a morphism $U \to F$ which is
surjective, flat, and locally of finite presentation.
\end{enumerate}
\end{theorem}

\begin{proof}
Trivial observations: (6) is a special case of (5) and
(4) is a special case of (3).
We first prove that cases (5) and (3) reduce to case (1).
Namely, by bootstrapping the diagonal
Lemma \ref{lemma-bootstrap-diagonal}
we see that (3) implies (5). In case (5) we set $R = U \times_F U$ which
is an algebraic space by assumption. Moreover, by assumption both
projections $s, t : R \to U$ are surjective, flat and locally of
finite presentation. The map $j : R \to U \times_S U$ is clearly an
equivalence relation. By
Lemma \ref{lemma-surjective-flat-locally-finite-presentation}
the map $U \to F$ is a surjection of sheaves. Thus $F = U/R$
which reduces us to case (1).

\medskip\noindent
Next, we show that (1) reduces to (2).
Namely, let $(U, R, s, t, c)$ be a groupoid in algebraic spaces
over $S$ such that $s, t$ are flat and locally of finite presentation, and
$j = (t, s) : R \to U \times_S U$ is an equivalence relation.
Choose a scheme $U'$ and a surjective \'etale morphism $U' \to U$.
Let $R' = R|_{U'}$ be the restriction of $R$ to $U'$. By
Groupoids in Spaces,
Lemma \ref{spaces-groupoids-lemma-quotient-pre-equivalence-relation-restrict}
we see that $U/R = U'/R'$. Since $s', t' : R' \to U'$ are also
flat and locally of finite presentation (see
More on Groupoids in Spaces,
Lemma \ref{spaces-more-groupoids-lemma-restrict-preserves-type})
this reduces us to the case where $U$ is a scheme.
As $j$ is an equivalence relation we see that $j$ is a monomorphism.
As $s : R \to U$ is locally of finite presentation we see that
$j : R \to U \times_S U$ is locally of finite type, see
Morphisms of Spaces, Lemma \ref{spaces-morphisms-lemma-permanence-finite-type}.
By
Morphisms of Spaces, Lemma
\ref{spaces-morphisms-lemma-monomorphism-loc-finite-type-loc-quasi-finite}
we see that $j$ is locally quasi-finite and separated.
Hence if $U$ is a scheme, then $R$ is a scheme by
Morphisms of Spaces, Proposition
\ref{spaces-morphisms-proposition-locally-quasi-finite-separated-over-scheme}.
Thus we reduce to proving the theorem in case (2).

\medskip\noindent
Assume $F = U/R$ where $(U, R, s, t, c)$ is a groupoid scheme
over $S$ such that $s, t$ are flat and locally of finite presentation, and
$j = (t, s) : R \to U \times_S U$ is an equivalence relation. By
Lemma \ref{lemma-slice-equivalence-relation}
we reduce to that case where $s, t$ are flat,
locally of finite presentation, and locally quasi-finite.
Let $U = \bigcup_{i \in I} U_i$ be an affine open covering
(with index set $I$ of cardinality $\leq$ than the size of $U$ to avoid
set theoretic problems later -- most readers can safely ignore this remark).
Let $(U_i, R_i, s_i, t_i, c_i)$ be the restriction of $R$
to $U_i$. It is clear that $s_i, t_i$ are still flat, locally of finite
presentation, and locally quasi-finite as $R_i$ is the open subscheme
$s^{-1}(U_i) \cap t^{-1}(U_i)$ of $R$
and $s_i, t_i$ are the restrictions of $s, t$ to this open. By
Lemma \ref{lemma-better-finding-opens}
(or the simpler
Spaces, Lemma \ref{spaces-lemma-finding-opens})
the map $U_i/R_i \to U/R$ is representable by open immersions.
Hence if we can show that $F_i = U_i/R_i$ is an algebraic space, then
$\coprod_{i \in I} F_i$ is an algebraic space by
Spaces, Lemma \ref{spaces-lemma-coproduct-algebraic-spaces}.
As $U = \bigcup U_i$ is an open covering it is clear that
$\coprod F_i \to F$ is surjective. Thus
it follows that $U/R$ is an algebraic space, by
Spaces, Lemma \ref{spaces-lemma-glueing-algebraic-spaces}.
In this way we reduce to the case where $U$ is affine and $s, t$ are flat,
locally of finite presentation, and locally quasi-finite and
$j$ is an equivalence.

\medskip\noindent
Assume $(U, R, s, t, c)$ is a groupoid scheme over $S$,
with $U$ affine, such that $s, t$ are flat, locally of finite presentation,
and locally quasi-finite, and $j$ is an equivalence relation.
Choose $u \in U$. We apply
More on Groupoids in Spaces,
Lemma \ref{spaces-more-groupoids-lemma-quasi-splitting-affine-scheme}
to $u \in U, R, s, t, c$. We obtain an affine scheme $U'$, an \'etale
morphism $g : U' \to U$, a point $u' \in U'$ with $\kappa(u) = \kappa(u')$
such that the restriction $R' = R|_{U'}$ is quasi-split over $u'$.
Note that the image $g(U')$ is open as $g$ is \'etale and contains $u$.
Hence, repeatedly applying the lemma, we can find finitely many
points $u_i \in U$, $i = 1, \ldots, n$,
affine schemes $U'_i$, \'etale morphisms $g_i : U_i' \to U$, points
$u'_i \in U'_i$ with $g(u'_i) = u_i$ such that (a) each
restriction $R'_i$ is quasi-split over some point in $U'_i$ and
(b) $U = \bigcup_{i = 1, \ldots, n} g_i(U'_i)$.
Now we rerun the last part of the argument in the preceding paragraph:
Using
Lemma \ref{lemma-better-finding-opens}
(or the simpler
Spaces, Lemma \ref{spaces-lemma-finding-opens})
the map $U'_i/R'_i \to U/R$ is representable by open immersions.
If we can show that $F_i = U'_i/R'_i$ is an algebraic space, then
$\coprod_{i \in I} F_i$ is an algebraic space by
Spaces, Lemma \ref{spaces-lemma-coproduct-algebraic-spaces}.
As $\{g_i : U'_i \to U\}$ is an \'etale covering
it is clear that $\coprod F_i \to F$ is surjective. Thus
it follows that $U/R$ is an algebraic space, by
Spaces, Lemma \ref{spaces-lemma-glueing-algebraic-spaces}.
In this way we reduce to the case where $U$ is affine and $s, t$ are flat,
locally of finite presentation, and locally quasi-finite,
$j$ is an equivalence, and $R$ is quasi-split over $u$ for some
$u \in U$.

\medskip\noindent
Assume $(U, R, s, t, c)$ is a groupoid scheme over $S$,
with $U$ affine, $u \in U$ such that $s, t$ are flat, locally
of finite presentation, and locally quasi-finite and
$j = (t, s) : R \to U \times_S U$ is an equivalence relation
and $R$ is quasi-split over $u$. Let $P \subset R$ be a quasi-splitting
of $R$ over $u$. By
Lemma \ref{lemma-divide-subgroupoid}
we see that $(U, R, s, t, c)$ is the restriction of a groupoid
$(\overline{U}, \overline{R}, \overline{s}, \overline{t}, \overline{c})$
by a surjective finite locally free morphism $U \to \overline{U}$ such that
$P = U \times_{\overline{U}} U$. Note that $s$ admits a factorization
$$
R = U \times_{\overline{U}, \overline{t}} \overline{R}
\times_{\overline{s}, \overline{U}} U
\xrightarrow{\text{pr}_{23}}
\overline{R} \times_{\overline{s}, \overline{U}} U
\xrightarrow{\text{pr}_2} U
$$
The map $\text{pr}_2$ is the base change of $\overline{s}$, and
the map $\text{pr}_{23}$ is a base change of the surjective finite locally
free map $U \to \overline{U}$. Since $s$ is flat, locally
of finite presentation, and locally quasi-finite and since $\text{pr}_{23}$
is surjective finite locally free (as a base change of such), we
conclude that $\text{pr}_2$ is flat, locally
of finite presentation, and locally quasi-finite by
Descent, Lemmas
\ref{descent-lemma-flat-fpqc-local-source} and
\ref{descent-lemma-locally-finite-presentation-fppf-local-source} and
Morphisms, Lemma \ref{morphisms-lemma-quasi-finite-local-source}.
Since $\text{pr}_2$ is the base change of the morphism
$\overline{s}$ by $U \to \overline{U}$ and $\{U \to \overline{U}\}$
is an fppf covering we conclude $\overline{s}$ is
flat, locally of finite presentation, and locally quasi-finite, see
Descent, Lemmas \ref{descent-lemma-descending-property-flat},
\ref{descent-lemma-descending-property-locally-finite-presentation}, and
\ref{descent-lemma-descending-property-quasi-finite}. The same goes
for $\overline{t}$. Consider the commutative diagram
$$
\xymatrix{
U \times_{\overline{U}} U \ar@{=}[r] \ar[rd] & P \ar[r] \ar[d] & R \ar[d] \\
& \overline{U} \ar[r]^{\overline{e}} & \overline{R}
}
$$
It is a general fact about restrictions that the outer four corners
form a cartesian diagram. By the equality we see the inner square is
cartesian. Since $P$ is open in $R$ (by definition of a quasi-splitting)
we conclude that $\overline{e}$ is an open immersion by
Descent, Lemma \ref{descent-lemma-descending-property-open-immersion}.
An application of
Groupoids,
Lemma \ref{groupoids-lemma-quotient-pre-equivalence-relation-restrict}
shows that $U/R = \overline{U}/\overline{R}$. Hence we have reduced to
the case where $(U, R, s, t, c)$ is a groupoid scheme over $S$,
with $U$ affine, $u \in U$ such that $s, t$ are flat, locally
of finite presentation, and locally quasi-finite and
$j = (t, s) : R \to U \times_S U$ is an equivalence relation
and $e : U \to R$ is an open immersion!

\medskip\noindent
But of course, if $e$ is an open immersion and
$s, t$ are flat and locally of finite presentation
then the morphisms $t, s$ are \'etale.
For example you can see this by applying
More on Groupoids, Lemma \ref{more-groupoids-lemma-sheaf-differentials}
which shows that $\Omega_{R/U} = 0$ which in turn implies
that $s, t : R \to U$ is G-unramified (see
Morphisms, Lemma \ref{morphisms-lemma-unramified-omega-zero}),
which in turn implies that $s, t$ are \'etale (see
Morphisms, Lemma \ref{morphisms-lemma-flat-unramified-etale}).
And if $s, t$ are \'etale then finally $U/R$ is an algebraic
space by
Spaces, Theorem \ref{spaces-theorem-presentation}.
\end{proof}

\end{document}
